\section{Defense Mechanisms}
\label{sec:defenses}

We evaluate five defenses under a \emph{pre-ingestion filtering} model:
before any new entry is stored in memory, the defense is applied; flagged
entries are blocked.  Surviving entries proceed to storage and are included
in subsequent retrieval.  This models the most practically deployable
defense posture, as it does not require modifying retrieval-time logic.

\subsection{Watermark Defense}

Content provenance is established at write time using the unigram watermark
of \citet{zhao2024watermark}.  The vocabulary is partitioned into green and
red lists based on a 256-bit secret key, and content is biased toward
green-list characters (proportionally controlled by $\gamma=0.25$,
$\delta=2.0$).  At read time, the z-score
\begin{equation}
  z = \frac{g - \gamma n}{\sqrt{\gamma(1-\gamma)n}}
\end{equation}
is computed, where $g$ is the count of green-list characters and $n$ is the
content length.  Entries with $z < z_{\mathrm{thr}} = 4.0$ are flagged as
potentially unsigned (adversarial).  This is a character-level adaptation of
the original token-level scheme, suitable for memory entries that do not
pass through a generative model.

\subsection{Input Validation Defense}

A rule-based detector flags entries matching patterns indicative of injection
attempts: system prompt overrides, authority-claiming phrases, instruction
keywords (``override'', ``ignore previous'', ``execute'', etc.).  Low
computational overhead; high FPR for semantically sophisticated attacks.

\subsection{Proactive Defense}

Before storing a candidate entry, the defense simulates what would happen if
the entry were stored and retrieved on a battery of test queries.  Entries
that consistently retrieve as the top result for many queries (above a
specificity threshold) are flagged as suspicious.  This is analogous to
\citet{gao2019strip}'s perturbation-based detection.

\subsection{Composite Defense}

A weighted ensemble of Watermark, Validation, and Proactive detectors.
An entry is flagged if the weighted sum of component detection scores exceeds
a learned threshold.  Achieves the highest $\tpr$ and lowest $\fpr$ of the
single-round defenses.

\subsection{Semantic Anomaly Detection (\sad{})}
\label{sec:sad}

\paragraph{Motivation.}
\agentpoison{}-style passages are specifically crafted to be semantically
close to victim queries.  This closeness, necessary for the attack to succeed,
is a detectable signal: \emph{legitimate memory entries from benign sources
are not expected to be highly similar to observed victim queries}.

\paragraph{Algorithm.}
\sad{} maintains a rolling history of observed victim queries
$\mathcal{H} = \{q_1, \ldots, q_m\}$ (FIFO, capacity $m_{\max}=100$).
After calibrating on a benign reference corpus, it computes for each
candidate entry $c$:
\begin{equation}
  s(c) = \max_{q \in \mathcal{H}} \cos(E(c), E(q)).
  \label{eq:sad_score}
\end{equation}
The calibration step estimates $\mu$ and $\sigma$ of $s(\cdot)$ over benign
entries and observed queries.  An entry is flagged as anomalous if:
\begin{equation}
  s(c) > \mu + k \cdot \sigma, \qquad k = 2.0\;\text{(default)}.
  \label{eq:sad_threshold}
\end{equation}

\paragraph{Properties.}
\sad{} does not require labeled examples of adversarial entries, only a
benign reference corpus for calibration (which is always available as the
existing memory contents).  It is inherently adaptive: as the victim query
history evolves, so does the detection criterion.  Batch encoding via
\texttt{detect\_batch()} amortizes encoder inference over all candidate
entries in a single forward pass.

\paragraph{Limitations.}
\sad{} is less effective against \injecmem{}, whose broad-anchor passages are
specifically designed \emph{not} to be suspiciously similar to any particular
query.  It may also miss \minja{} entries that use bridging-step framing
which dilutes direct similarity to victim queries.
